\documentclass[10pt]{article}

\usepackage[utf8]{inputenc}

\usepackage[T1]{fontenc}

\usepackage{amsmath}

\usepackage{amsfonts}

\usepackage{amssymb}

\usepackage[version=4]{mhchem}

\usepackage{stmaryrd}

\usepackage{graphicx}

\usepackage[export]{adjustbox}

\graphicspath{ {./images/} }



\begin{document}

\begin{enumerate}

  \setcounter{enumi}{1}

  \item Л и н е й в о е с в о й с т в о: для любых постоянных $\alpha$ и $\beta$ из интегрируемости на $[a, b]$ каждой из функций $f(x)$ и $g(x)$ следуют интегрируемость на этом отрезке функции $[\alpha f(x)+\beta g(x)]$ и равенство

\end{enumerate}



\[

\int_{a}^{b}[\alpha f(x)+\beta g(x)] d x=\alpha \int_{a}^{b} f(x) d x+\beta \int_{a}^{b} g(x) d x .

\]



\begin{enumerate}

  \setcounter{enumi}{2}

  \item Из интегрируемости на отрезке $[a, b]$ каждой из функций $f(x)$ и $g(x)$ следует интегрируемость на этом отрезке произведения $f(x) g(x)$.



  \item А д д и т и в н о с т Б: из интегрируемости функции $f(x)$ на каждом из отрезков $[a, c]$ и $[c, b]$ следуют интегрируемость $f(x)$ на отрезке $[a, b]$ в равенство



\end{enumerate}



\[

\int_{a}^{b} f(x) d x=\int_{a}^{c} f(x) d x+\int_{c}^{b} f(x) d x .

\]



\begin{enumerate}

  \setcounter{enumi}{4}

  \item Если функции $f(x)$ и $g(x)$ интегрируемы на $[a, b]$ и если $f(x) \geqslant g(x)$ всюду на этом отрезке, то

\end{enumerate}



\[

\int_{a}^{b} f(x) d x \geqslant \int_{a}^{b} g(x) d x .

\]



\begin{enumerate}

  \setcounter{enumi}{5}

  \item Из интегрируемости на $[a, b]$ функции $f(x)$ следуют интегрируемость на этом отрезке функции $|f(x)|$ и справедливость оценки

\end{enumerate}



\[

\left|\int_{a}^{b} f(x) d x\right| \leqslant \int_{a}^{b}|f(x)| d x

\]



\begin{enumerate}

  \setcounter{enumi}{6}

  \item Ф о р м ул а с ред не го значения: если функции $f(x)$ и $g(x)$ интегрируемы на $[a, b]$, функция $\boldsymbol{g}(x)$ неотрицательна или неположительна всюду на этом отрезке, а $M$ и $m$ - точные верхняя и нижняя грани $f(x)$ на $[a, b]$, то найдется число $\mu$ из отрезка $m \leqslant \mu \leqslant M$ такое, что справедлива формула

\end{enumerate}



\[

\int_{a}^{b} f(x) g(x) d x=\mu \int_{a}^{b} g(x) d x .

\]



Если, кроме того, функция $f(x)$ непрерывна на $[a, b]$, то на этом отрезке найдется точка $\xi$ такая, что в формуле (3)



\[

\mu=f(\xi) .

\]



\begin{enumerate}

  \setcounter{enumi}{7}

  \item В торая фо рмула с реднего значен и я (фо р м у л а Бон н е): если функция $f(x)$ интегрируема на $[a, b]$, а функция $g(x)$ монотонна на этом отрезке, то найдется точка $\xi$ на $[a, b]$ такая, что справедлива формула

\end{enumerate}



\[

\int_{a}^{b} f(x) g(x) d x=g(a) \int_{a}^{\xi} f(x) d x+g(b) \int_{\xi}^{b} f(x) d x .

\]



Jum.: [1] R i e m a n n B., "Göttinger Akad. Abhandl.», 1868, Вd 13; [2] И ль и н В. А., П о з н я к Ә. Г., Основы математического анализа, 3 изд., ч. 1, М., 1971, 2 изд., ч. 2, М , 1980; [3] К у д р я в ц е в Л. Д., курс математического анализа, ческого анализа, 2 изд., т. $1-2$, М., 1975 . М., Кур. м. В. Ильин.



\includegraphics[max width=\textwidth, center]{2023_07_08_b478b332a06e68c92516g-1}

ме т од,-- метод решения Гурса задачи и Коши задачи для линейных гиперболич. типа уравнений 2-го порядка с двумя независимыми переменными



\[

\begin{gathered}

L u \equiv u_{x y}+a(x, y) u_{x}+b(x, y) u_{y}+c(x, y) u= \\

=f(x, y) .

\end{gathered}

\]



В Р. м. фундаментальную роль играет ф у н к ц и я Р и м а н а $R=R(x, y ; \xi, \eta)$ к-рая ири определенных предположениях относительно заданных функциї $a, b$, $c$ и $f$ однозначно определяется как решение специальвой задачи Гурса:



\[

\begin{aligned}

& R(\xi, y ; \xi, \eta)=\exp \int_{\eta}^{y} a(\xi, t) d t ; \\

& R(x, \eta ; \xi, \eta)=\exp \int_{\xi}^{x} b(t, \eta) d t

\end{aligned}

\]



для сопряженного уравнения



\[

L * R \equiv R_{x y}-\frac{\partial}{\partial x}(a, R)-\frac{\partial}{\partial y}(b R)+c R=0 .

\]



Функция $R$ по переменным $\xi, \eta$ является решением однородного уравнения



\[

R_{\xi \eta}+a(\xi, \eta) R_{\xi}+b(\xi, \eta) R_{\eta}+c(\xi, \eta) R=0 .

\]



При $a=b=0, c=$ const функция $R=J_{0}(\sqrt{4 c(x-\xi)(y-\eta)})$, где $J_{0}(z)$ - Функция Бесселя порядка нуль.



Функцию Римана можно определить как решение нагруженного интегрального уравнения Вольтерра:



\[

\begin{gathered}

R(x, y ; \xi, \eta)-\int_{\eta}^{y} a(x, \tau) R(x, \tau ; \xi, \eta) d \tau- \\

-\int_{\xi}^{x} b(t, y) R(t, y ; \xi, \eta) d t+ \\

+\int_{\xi}^{x} d t \int_{\eta}^{y} c(t, \tau) R(t, \tau ; \xi, \eta) d \tau=1 .

\end{gathered}

\]



Р. м. решения задачи Гурса реализуется следующим образом: для любой дифференцируемой до соответствующего порядка функции $u=u(x, y)$ имеет место тождество



\[

\begin{aligned}

& \frac{\partial^{2}}{\partial x \partial y}[u R(x, y ; \xi, \eta)]-R(x, y ; \xi, \eta) L u= \\

& =\frac{\partial}{\partial x}\left[u\left(\frac{\partial R}{\partial y}-a R\right)\right]+\frac{\partial}{\partial y}\left[u\left(\frac{\partial R}{\partial x}-b R\right)\right],

\end{aligned}

\]



из к-рого интегрированием по частям получается, что любое решение $u$ уравнения (1) представляет собой решение нагруженного интегрального уравнения:



\[

\begin{gathered}

u(x, y)=R\left(x, y_{0} ; x, y\right) u\left(x, y_{0}\right)+ \\

+R\left(x_{0}, y ; \xi, y\right) u\left(x_{0}, y\right)-R\left(x_{0}, y_{0} ; x, y\right)+ \\

+\int_{x_{0}}^{x}\left[b\left(t, y_{0}\right) R\left(t, y_{0} ; x, y\right)-\frac{\partial R\left(t, y_{0} ; x, y\right)}{\partial t}\right] \times \\

\times u\left(t, y_{0}\right) d t+\int_{y_{0}}^{y}\left[a\left(x_{0}, \tau\right) R\left(x_{0}, \tau ; x, y\right)-\right. \\

\left.-\frac{\partial R\left(x_{0}, \tau ; x, y\right)}{\partial \tau}\right] u\left(x_{0}, \tau\right) d \tau+ \\

+\int_{x_{0}}^{x} d t \int_{y_{0}}^{y} R(t, \tau ; x, y) f(x, \tau) d \tau, x>x_{0}, y>y_{0} .

\end{gathered}

\]



Из (3) непосредственно вытекает корректность задачи Гурса



\[

u\left(x, y_{0}\right)=\varphi(x), u\left(x_{0}, y_{0}\right)=\psi(y), \varphi\left(x_{0}\right)=\psi\left(y_{0}\right)

\]



для уравнения (1).



Р. м. прпводит решение задачи Коши для уравнения (1) с начальными данными на любой гладкой нехарактеристич. кривой к нахождению функции Римана и дает возможность в квадратурах выписать решение этой задачи.



Р. м. обобщен на широкий класс линейных гиперболич. уравнений и систем.



Для случая гиперболич. типа системы линейных уравнений с частными производными 2-го порядка



$u_{x x}-u_{y y}+a(x, y) u_{x}+b(x, y) u_{y}+c(x, y) u=f(x, y)$, где $a, b, c$ - заданные действительные квадратные симметрич. матрицы порядка $m, f=\left(f_{1}, \ldots, f_{m}\right)$ - заданныӥ, a $u=\left(u_{1}, \ldots, u_{m}\right)$ - искомый векторы, м а т р и ц а Р и а в а однозначно определяется как решіение системы нагруженных интегральных уравнений Вольтерра вида (2), в правой части к-рой стоит единичная матрица $I$ порядка $m$.



В. Вольтерра (V. V olterra) впервые обобщил Р. м. на волновое уравнение



\[

u_{x x}+u_{y y}-u_{t t}=f(x, y, t)

\]





\end{document}